\documentclass[11pt,a4paper]{article}
\pdfextension mapfile{+cm.map}
\pdfextension mapfile{+cmextra.map}
\pdfextension mapfile{+symbols.map}
\pdfextension mapfile{+euler.map}
\pdfextension mapfile{+latxfont.map}
\pdfextension mapfile{+rsfs.map}
\usepackage[margin=25mm,headheight=15pt]{geometry}
\usepackage{fontspec}
\setmainfont{Latin Modern Roman}
\setsansfont{Latin Modern Sans}
\setmonofont{Latin Modern Mono}[Scale=MatchLowercase]
\usepackage{amsmath,amssymb,amsthm,mathtools,mathrsfs}
\usepackage{booktabs,longtable,array,tabularx,multirow}
\usepackage{graphicx,xcolor,microtype,fancyhdr}
\usepackage[unicode,hidelinks]{hyperref}
\usepackage{xurl,enumitem,listings}
\makeatletter\renewcommand{\@pnumwidth}{2.5em}\makeatother
\lstset{basicstyle=\ttfamily\scriptsize,breaklines=true,columns=fullflexible,keepspaces=true}
\allowdisplaybreaks
\setcounter{MaxMatrixCols}{64}
\emergencystretch=2em
\setlength{\parindent}{1.3em}
\setlength{\parskip}{2pt}
\newcommand{\Z}{\mathbb Z}
\newcommand{\Q}{\mathbb Q}
\newcommand{\R}{\mathbb R}
\newcommand{\C}{\mathbb C}
\newcommand{\Pp}{\mathbb P}
\newcommand{\OO}{\mathcal O}
\newcommand{\HH}{\mathcal H}
\newcommand{\LL}{\mathcal L}
\newcommand{\cT}{\mathcal T}
\newcommand{\id}{\mathrm{id}}
\newcommand{\rank}{\operatorname{rank}}
\newcommand{\im}{\operatorname{im}}
\newcommand{\coker}{\operatorname{coker}}
\newcommand{\Aut}{\operatorname{Aut}}
\newcommand{\Spec}{\operatorname{Spec}}
\newcommand{\diag}{\operatorname{diag}}
\newcommand{\ord}{\operatorname{ord}}
\newcommand{\Hom}{\operatorname{Hom}}
\newcommand{\tr}{\operatorname{tr}}
\newcommand{\sgn}{\operatorname{sgn}}
\newcommand{\red}{\mathrm{red}}
\newcommand{\vir}{\mathrm{vir}}
\newcommand{\HP}{\mathbb H\mathbb P}
\newenvironment{statement}[2]{\par\medskip\noindent\textbf{Statement #1 (#2).}\label{stmt:#1}\quad\ignorespaces}{\par\medskip}
\pagestyle{fancy}\fancyhf{}
\fancyhead[L]{\small S6\_06\_BLIND}\fancyhead[R]{\small Independent derivations}
\fancyfoot[C]{\thepage}
\title{S6\_06\_BLIND\\Independent derivations from the supplied constructions}
\author{AI-produced by GPT-6 Astra\\OpenAI, ChatGPT Pro mode}
\date{18 November 2026\\Working manuscript; unresolved obligations remain explicit}
\begin{document}
\maketitle
\noindent We use the complex orientation on every complex manifold and the positive, counterclockwise meridian in each holomorphic transverse coordinate. Put $B=\Pp^1_t$, $p_1=0$, $p_2=1$, $p_0=\infty$, and let $S_j$ denote the reduced multiple fibre. In the $(3,4)$ construction the full finite fibres are $3S_1$ and $4S_2$. The ordered regular-fibre lattice and periods are
\[
 \Lambda=\Z\langle\widehat\gamma,\widehat u,\widehat w,\widehat\delta\rangle,
 \qquad\Pi_c=\begin{pmatrix}6\mu&\tau&1&0\\ \beta^{\rm part}+c&\mu&0&1\end{pmatrix}.
\]
We use $w=e^{-\pi i c}$ and distinguish this parameter from the period $\widehat w$. At the cusp $q=t^{-1}$, $s=\log q/(2\pi i)$; after $w=r^6$ put $v=\log r/(2\pi i)$ and $c=-12v$. Positive $q,r$ meridians increase $s,v$ by one. A period-column comparison $\Pi'=\Pi K$ induces $K^{-1}$ on fixed vector coordinates.

We identify $\Theta_7=\Z/28$ by $[\Sigma]\mapsto28\mu(\Sigma)$, with outward-normal-first boundary orientation and $\mu=(p_1^2-4\sigma)/896$ for a spin bound. The positive generator has $\mu=1/28$. The forward mapping torus uses $(x,1)\sim(\Phi(x),0)$; the reverse convention used for the six-dimensional return and the boundary-orientation comparison are recorded explicitly in P3K.7. P3F.22 proves the transfer with these orientations.

The comparison with Engel's construction uses the holomorphic characters $Z=e^{2\pi i\zeta_1}$, $W=e^{-2\pi i\zeta_2}$. The Tate period is the $\widehat u$ column; the lifted translation is the $\widehat\gamma$ column; fibrewise $M^\times$ scaling is the $W$ direction, with primitive circle $\widehat\delta$ subject to this exponential sign. Engel's $h=t_{6P}$ and $\phi:h^*M\to M$ give the generator covering $h^{-1}$. Its multiplier is the $\widehat\gamma$ multiplier on $W$. Thus, after fixing the reference $\phi$, Engel's scalar parameter is $u=\kappa w^2$ for a nonzero constant $\kappa$ fixed by that reference: both multiply the $\widehat\gamma$ character by $e^{-2\pi ic}$. We retain the unspecified additive normalization of $\beta^{\rm part}$. In Part~1, $A,B$ act on cohomology and $M_1=A^t,M_2=B^t$ on homology.
\tableofcontents\clearpage
\IfFileExists{notes/part1.tex}{\input{notes/part1.tex}}{The current proofs and exact computations are preserved in \texttt{notes/part1.md} and \texttt{code/part1/}.}
\IfFileExists{notes/part1_analytic.tex}{\input{notes/part1_analytic.tex}}{}
\IfFileExists{notes/part1_endpoints.tex}{\input{notes/part1_endpoints.tex}}{}
\input{notes/part1_arithmetic.tex}
\section{The corner and the parameter origin}
\IfFileExists{notes/part2.tex}{\input{notes/part2.tex}}{The current proofs are preserved in the Part~2 notes.}
\section{Smooth monodromy and the page}
\IfFileExists{notes/part3.tex}{\input{notes/part3.tex}}{The current proofs and unresolved geometric comparisons are preserved in \path{notes/part3_family.md} and \path{code/part3/}.}
\input{notes/part3_characteristics.tex}
\section{The threefold checks}
\IfFileExists{notes/part4.tex}{\input{notes/part4.tex}}{The current proof is preserved in \texttt{notes/part4.md}.}
\IfFileExists{notes/part4_unibranch.tex}{\input{notes/part4_unibranch.tex}}{}
\section{The blind control}
\IfFileExists{notes/part5.tex}{\input{notes/part5.tex}}{The independent proof is preserved in \texttt{notes/part5\_blind.md}. The comparison paper has not been opened.}
\IfFileExists{notes/part5_comparison.tex}{\input{notes/part5_comparison.tex}}{}
\nocite{*}
\bibliographystyle{alpha}\bibliography{S6_06_BLIND_references}
\end{document}
